\section*{平面机构自由度计算}

\subsection*{核心计算公式}
平面机构自由度采用葛氏（契比雪夫）公式：
\begin{equation}
F = 3n - 2p_l - p_h
\end{equation}
其中：
\begin{itemize}
    \item $F$：机构自由度（Degrees of Freedom）；
    \item $n$：活动构件个数（Number of moving links）；
    \item $p_l$：低副个数（Number of lower pairs），包括转动副和移动副；
    \item $p_h$：高副个数（Number of higher pairs），包括齿轮副和凸轮副。
\end{itemize}

\subsection*{边界与约束条件（测试关键点）}

\subsubsection*{1. 复合铰链}
$m$ 个构件在同一处构成轴线重合的转动副，铰链数为 $m - 1$。

\subsubsection*{2. 局部自由度}
构件中不影响整个机构运动规律的独立运动（如凸轮机构中的滚子自转），计算时需将其排除，即减去该局部自由度。

\subsubsection*{3. 虚约束}
对机构的运动不起实际限制作用的重复约束（如平行双曲柄机构中的对称连杆），计算时应予以排除。

\subsection*{机构具有确定运动的条件}
机构的自由度数 $F$ 必须大于 0，且必须等于机构的原动件（输入件）数。